\begin{exercise}{18}
Show that a natural transformation between two functors
$f,g\colon\mathcal C\to\mathcal D$ induces a homotopy between
$\mathrm N(f),\mathrm N(g)\colon\mathrm N(\mathcal C)\to
\mathrm N(\mathcal D)$. Use this result to show that conjugation
with an element determines a self map of $BG$ which is homotopic
to the identity. What does conjugation induce on $\pi_1(BG)$?
Why does this not show that every group is abelian?

\begin{proof}
Let $\alpha\colon f\Rightarrow g$ be the natural transformation.
Define $A\colon\mathcal C\times[1]\to\mathcal D$ by
$A(c,0)=f(c)$ and $A(c,1)=g(c)$. On arrows at level $0$ or $1$,
use $f$ or $g$, respectively; for a mixed arrow $(a,0\to1)$ with
$a\colon c\to c'$, put
\[
 A(a,0\to1)=\alpha_{c'}f(a)=g(a)\alpha_c.
\]
Naturality makes $A$ a functor. The nerve preserves products by
its definition, and $\mathrm N([1])\cong\Delta^1$ by
Lemma~1.1.44. Thus $\mathrm N(A)$ gives
\[
 \mathrm N(\mathcal C)\times\Delta^1\longrightarrow
 \mathrm N(\mathcal D)
\]
with endpoints $\mathrm N(f)$ and $\mathrm N(g)$, the required
simplicial homotopy. Realization gives an ordinary homotopy by
Remark~1.2.60.

Regard $G$ as a category with one object $*$ and put
$c_a(h)=aha^{-1}$ for $a\in G$.
The arrow $a$ is a natural transformation
$\mathrm{id}\Rightarrow c_a$, since $c_a(h)a=ah$.
Hence the induced map $Bc_a\colon BG\to BG$ is homotopic to the
identity.

The CW structure obtained from Exercise~\ref{ex:15}
(Corollary~1.1.41) presents $\pi_1(BG,*)$ with
edge generators $[h]$ for $h\in G$ and relations
$[1]=1$ and $[k][h]=[kh]$ from the $2$-simplices.
Thus $[h]\mapsto h$ identifies $\pi_1(BG,*)$ with $G$, and
$Bc_a$ induces the inner automorphism $h\mapsto aha^{-1}$.
The homotopy above moves the basepoint along the loop represented
by $a$. It is a free homotopy, which permits conjugate induced maps
on the fundamental group; it does not require the conjugation
automorphism to be the identity. Thus it gives no conclusion that
$G$ is abelian.
\end{proof}
\end{exercise}
